\section*{الفصل \ref{ch:b3:climate-models} --- النماذج المناخية}

\begin{solution}{exo:b3:climate-models:1}
$A = 14.0 - 1.5 = \num{12.5}$ و$B = 0.14 - 0.05 = \num{0.09}$، بوحدة
\unit{kcal.cm^{-2}} في الشهر. والشهر يساوي \qty{2.63e6}{s}، فإنّ
\qty{1}{kcal.cm^{-2}} في الشهر تساوي $\num{4.184e7}/\num{2.63e6} =
\qty{15.91}{W/m^2}$: $A = \qty{199}{W/m^2}$، $B =
\qty{1.43}{W.m^{-2}.\celsius^{-1}}$. عند \qty{15}{\celsius}، $199 + 21.5 =
\qty{220}{W/m^2}$.
\end{solution}

\begin{solution}{exo:b3:climate-models:2}
$(238.0 - 211.9)/1.816 = \qty{14.4}{\celsius}$؛ مع 0.31، $(234.6 -
211.9)/1.816 = \qty{12.5}{\celsius}$. كلّ 0.01 من الألبيدو يكلّف
\qty{3.4}{W/m^2} و$3.4/1.816 = \qty{1.9}{\celsius}$.
\end{solution}

\begin{solution}{exo:b3:climate-models:3}
$C = 1025\times3992\times60 = \qty{2.46e8}{J.m^{-2}.K^{-1}}$؛ بقسمتها على
\qty{3.156e7}{s} من السنة، \qty{7.78}{W.yr.m^{-2}.K^{-1}}. ثم $\tau
= 7.78/1.816 = \num{4.3}\text{ سنة}$، و$7.78/1.33 = \num{5.8}\text{ سنة}$.
\end{solution}

\begin{solution}{exo:b3:climate-models:4}
\mbox{\textdir TLT ECS $= 3.93/1.33 = \qty{2.95}{\celsius}$}؛ \mbox{\textdir TLT TCR} $\approx 3.93/1.95 =
\qty{2.02}{\celsius}$؛ النسبة \num{0.68}.
\end{solution}

\begin{solution}{exo:b3:climate-models:5}
$Q = \qty{335}{W/m^2}$. كرة الثلج: $(335\times0.38 - 211.9)/1.816 =
\qty{-46.6}{\celsius}$؛ الحالة الدافئة: $(335\times0.70 - 211.9)/1.816 =
\qty{12.4}{\celsius}$، وما زالت أعلى من \qty{10}{\celsius}؛ المنحدر: $335(0.54 +
0.016\,T) = 211.9 + 1.816\,T$ فيعطي $T = \qty{8.7}{\celsius}$. $N' =
\qty{-1.816}{W.m^{-2}.\celsius^{-1}}$ عند الحالتين الأوليين، مستقرّة، و$-1.816
+ 335\times0.016 = \num{+3.54}$ عند المنحدر، غير مستقرّة.
\end{solution}

\begin{solution}{exo:b3:climate-models:6}
ينتهي الفرع الدافئ حيث تبلغ درجة حرارته \qty{0}{\celsius}:
$0.70\,Q = A$، $Q = \qty{302.7}{W/m^2}$، أي ثابت شمسي قدره
\qty{1211}{W/m^2}، أدنى بنسبة \qty{11.0}{\percent} من قيمته الحالية. تذوب كرة الثلج
حيث $0.38\,Q = A - 20B$، $Q = \qty{462.1}{W/m^2}$، $S_0 =
\qty{1848}{W/m^2}$. الحدّ الفاصل هو $B(T^* - T_F)/(Q_0(1 - \alpha_F))$،
الفرق بين درجة الحرارة الحالية $T^*$ ودرجة الحرارة التي عندها
يبدأ الجليد بالانتشار، محوّلًا إلى الإشماس: $1.816\times14.37/238 =
\qty{11.0}{\percent}$ هنا، و$1.816\times4.37/238 = \qty{3.3}{\percent}$ في
النصّ.
\end{solution}

\begin{solution}{exo:b3:climate-models:7}
(أ)~\qty{-45.5}{\celsius}. (ب)~$340\times0.38 + F = 211.9 - 18.2$: $F =
\qty{64.5}{W/m^2}$، أي 16 ضعف التأثير الإشعاعي لمضاعفة واحدة. (ج)~ما دام
التأثير قائمًا ينتقل الكوكب إلى الفرع الدافئ، $(238.0 + 64.5 -
211.9)/1.816 = \qty{50}{\celsius}$؛ فإذا أُزيل، يبرد الكوكب
على طول الفرع الدافئ حتى \qty{14.4}{\celsius}، وهذه الحالة قائمة عند الشمس الحالية:
فلا يعود إلى حالة كرة الثلج.
\end{solution}

\begin{solution}{exo:b3:climate-models:8}
$\Sigma = 1.80/7.0 + 0.70/90 = \qty{0.2649}{yr^{-1}}$، $\Pi = 1.10\times0.70/
(7.0\times90) = \qty{1.222e-3}{yr^{-2}}$. وتعطي الجذور $\tau_f =
\num{3.84}\text{ سنة}$ و$\tau_s = \num{213}\text{ سنة}$؛ أمّا التقريبات،
$7.0/1.80 = \num{3.89}\text{ سنة}$، و$90\times1.80/(1.10\times0.70) =
\num{210}\text{ سنة}$. $a_f = 0.597$ تمامًا، و$1.10/1.80 = 0.611$
تقريبًا. \mbox{\textdir TLT ECS $= 3.93/1.10 = \qty{3.57}{\celsius}$}؛ \mbox{\textdir TLT TCR} $\approx
3.93/1.80 = \qty{2.18}{\celsius}$.
\end{solution}

\begin{solution}{exo:b3:climate-models:9}
المقام $B + \beta = \num{5.555}$. عند الاستواء: $(340\times1.221\times0.70 -
211.9 + 3.74\times12.81)/5.555 = \qty{22.8}{\celsius}$؛ عند القطب:
$(340\times0.506\times0.38 - 211.9 + 47.9)/5.555 = \qty{-17.7}{\celsius}$.
يزيل النقل $3.74\times(22.8 - 12.8) = \qty{37}{W/m^2}$ عند
الاستواء، ويجلب $3.74\times30.6 = \qty{114}{W/m^2}$ إلى القطب.
\end{solution}

\begin{solution}{exo:b3:climate-models:10}
(أ)~$(1 - x^2)P_2' = 3x - 3x^3$، ومشتقّتها $3 - 9x^2 = -6P_2$.
(ب)~حدّ النقل هو $-\mathcal D T_2(-6P_2) = 6\mathcal DT_2P_2$.
الموازنة $Q(1 + s_2P_2)(1 - \alpha) - A - B(T_0 + T_2P_2) - 6\mathcal
DT_2P_2 = 0$ تتحقّق من أجل كلّ $x$ إذا انعدم جزؤها الثابت وجزؤها $P_2$
فنحصل على $T_0$ و$T_2$. (ج)~$T(0) - T(1) = T_2(-\tfrac12 - 1)
= -1.5\,T_2$. بدون النقل $T_2 = 340\times(-0.476)\times0.70/1.816 =
\qty{-62.4}{\celsius}$: بفرق قدره \qty{94}{\celsius}. (د)~$T_2 =
-55.3/1.5 = \qty{-36.9}{\celsius}$، فإنّ $B + 6\mathcal D = 113.3/36.9 =
3.07$ و$\mathcal D = \qty{0.21}{W.m^{-2}.\celsius^{-1}}$.
\end{solution}

\begin{solution}{exo:b3:climate-models:11}
(أ)~$r = \tfrac12\Sigma\big(1 \pm \sqrt{1 - 4\Pi/\Sigma^2}\big)$؛ مع
$\sqrt{1 - \epsilon} \approx 1 - \epsilon/2$، فتكون الجذور $\Sigma -
\Pi/\Sigma \approx \Sigma$ و$\Pi/\Sigma$. (ب)~عند $T_d = 0$، $C\,\mathrm
dT/\mathrm dt = kt - (\gamma - \lambda)T$، وحلّها انطلاقًا من السكون هو $T =
k(t - \tau_f)/(\gamma - \lambda) + k\tau_f\mathrm e^{-t/\tau_f}/(\gamma -
\lambda)$، $\tau_f = C/(\gamma - \lambda)$. من أجل $t \gg \tau_f$، $T \approx
kt/(\gamma - \lambda)$، وعند المضاعفة، $kt = F_{2\times}$: \mbox{\textdir TLT TCR}
$\approx F_{2\times}/(\gamma - \lambda)$. (ج)~$\Sigma = 1.95/8.1 +
0.62/110 = \qty{0.2464}{yr^{-1}}$، $\Pi = 1.33\times0.62/(8.1\times110) =
\qty{9.25e-4}{yr^{-2}}$؛ $\tau_f = \num{4.12}\text{ سنة}$ و$\tau_s =
\num{262}\text{ سنة}$ (تقريبًا $1/\Sigma = \num{4.06}\text{ سنة}$
و$\Sigma/\Pi = \num{266}\text{ سنة}$)؛ $4\Pi/\Sigma^2 = 0.061$.
\end{solution}

\begin{solution}{exo:b3:climate-models:12}
(أ)~المتوسط \qty{3.50}{\celsius}، والانحراف المعياري \qty{0.81}{\celsius}.
(ب)~\mbox{\textdir TLT ECS $= 1.88 + 1.86\,X$}؛ الترابط \num{0.992}؛ الانحراف المعياري
للباقي \qty{0.11}{\celsius}. (ج)~$1.88 + 1.86\times0.7 =
\qty{3.18}{\celsius}$، مع انتشار $\sqrt{(1.86\times0.1)^2 + 0.11^2} =
\qty{0.22}{\celsius}$، أي نحو ربع انتشار المجموعة.
(د)~لا توجد آلية فيزيائية تربط $X$ بالتغذية الراجعة؛ ومتغيّر $X$ وُجد
بتجربة مقادير رصدية كثيرة، بحيث قد يكون الترابط محض صدفة؛ وانحياز
مشترك بين النماذج الثمانية جميعها؛ ونماذج من العائلة نفسها تُعدّ
مستقلّة.
\end{solution}

\begin{solution}{pb:b3:climate-models:1}
\textbf{1.} $Q_0 = 1360/4 = \qty{340}{W/m^2}$: يعترض الكوكب
ضوء الشمس على قرص مساحته $\pi R^2$، ويوزّعه على كرة مساحتها
$4\pi R^2$.

\textbf{2.} $340\times0.70 = \qty{238.0}{W/m^2}$؛ $211.9 + 1.816\times15 =
\qty{239.1}{W/m^2}$.

\textbf{3.} $(238.0 - 211.9)/1.816 = \qty{14.4}{\celsius}$.

\textbf{4.} $(129.2 - 211.9)/1.816 = \qty{-45.5}{\celsius}$، أدنى من
\qty{-10}{\celsius}.

\textbf{5.} $340(0.54 + 0.016\,T) = 211.9 + 1.816\,T$: $T =
(183.6 - 211.9)/(1.816 - 5.44) = \qty{7.8}{\celsius}$.

\textbf{6.} $(238.0/\num{5.670e-8})^{1/4} = \qty{254.5}{K}$,
\qty{-18.6}{\celsius}: أي أدنى بمقدار \qty{33}{\celsius} من توازن الحالة الدافئة، ويعود ذلك إلى
تأثير الاحتباس الحراري للكوكب المدمج في $A$ و$B$.

\textbf{7.} $N' = -B = \qty{-1.816}{W.m^{-2}.\celsius^{-1}}$ عند
الكرة الثلجية وتوازن الحالة الدافئة، وكلاهما مستقرّ؛ $-1.816 + 5.44 =
\num{+3.62}$ عند التوازن الأوسط، وهو غير مستقرّ.

\textbf{8.} $\tau = 8.1/3.62 = \num{2.2}\text{ سنة}$.

\textbf{9.} $8.1/1.816 = \num{4.5}\text{ سنة}$.

\textbf{10.} $3.93/1.816 = \qty{2.2}{\celsius}$، أدنى من أفضل تقدير
مُقيَّم للمقدار \qty{3}{\celsius}، وأدنى من مداه \emph{المرجَّح}، من 2.5 إلى
\qty{4}{\celsius}: $B$ للنموذج، المُلاءَم عبر خطوط العرض، ليس
معامل التغذية الراجعة للمناخ ككلّ.

\textbf{11.} $\mathrm e^{-t/\tau} = 0.05$ عند $t = 3\tau =
\num{13}\text{ سنة}$. يُسخّن المناخ الحقيقي المحيط العميق أيضًا، الذي يُقاس زمنه
بالقرون.

\textbf{12.} عند $T_F = \qty{10}{\celsius}$؛ $Q = (211.9 + 18.2)/0.70 =
\qty{328.7}{W/m^2}$.

\textbf{13.} $S_0 = \qty{1315}{W/m^2}$، أدنى بنسبة \qty{3.3}{\percent} من قيمته اليوم.

\textbf{14.} بالنسبة إلى الكرة الثلجية، $(328.7\times0.38 - 211.9)/1.816 =
\qty{-47.9}{\celsius}$.

\textbf{15.} $Q = (211.9 - 18.2)/0.38 = \qty{509.8}{W/m^2}$، $S_0 =
\qty{2039}{W/m^2}$، أعلى بنسبة \qty{50}{\percent} من قيمته اليوم.

\textbf{16.} نزولًا على الفرع الدافئ حتى \qty{10}{\celsius} عند
\qty{1315}{W/m^2}، ثم انخفاض إلى \qty{-48}{\celsius}، ثم صعودًا على فرع الكرة الثلجية
حتى \qty{-45.5}{\celsius} عند \qty{1360}{W/m^2}، ويبقى الكوكب
متجمّدًا. ولا يعود إلى الفرع الدافئ إلّا فوق \qty{2039}{W/m^2}، حيث
يقفز مجددًا إلى الفرع الدافئ.

\textbf{17.} $\bar\alpha = 0.30\times0.9738 + 0.62\times0.0262 =
\num{0.3084}$; $\bar T = (340\times0.6916 - 211.9)/1.816 =
\qty{12.80}{\celsius}$.

\textbf{18.} $T_c = (340\times0.5643\times0.50 - 211.9 + 3.739\times12.80)/
5.555 = \qty{-12.26}{\celsius}$، أدفأ بمقدار \qty{3.5}{\celsius} من
إعادة التحليل عند \ang{72}.

\textbf{19.} $5.555\times(-12.26) + 211.9\times3.059 = \num{580.1}$، مع
$\beta/B = \num{2.059}$.

\textbf{20.} المقامات هي 1.7053، 1.7473، 1.7565، 1.7488،
و1.7362، و$Q/Q_0 = \num{1.0005}$، \num{0.9764}، \num{0.9713}،
\num{0.9756} و\num{0.9827}.

\textbf{21.} قرب \ang{47.5}، عند \num{0.971}. خطوط الجليد الواقعة نحو القطب من
الحدّ الأدنى مستقرّة، $\mathrm dQ/\mathrm dx_s > 0$؛ أما التي نحو خط الاستواء فهي
غير مستقرّة.

\textbf{22.} $Q/Q_0 = 1$ أيضًا بين \ang{30} و\ang{40}، قرب
\ang{30} (\ang{30.2} بالاستيفاء): خطّ جليد غير مستقرّ. عند \ang{0}،
$Q/Q_0 = 580.1/(0.6106 + 0.7824)/340 = \num{1.225}$: تبقى الكرة الثلجية قائمة
لأيّ شمس أسطع من شمس اليوم بنسبة تصل إلى \qty{22}{\percent}، لذا فهي
ممكنة اليوم.

\textbf{23.} \qty{2.9}{\percent}، مقابل \qty{3.3}{\percent} للنموذج
عديم البعد، و\qty{1.6}{\percent} لنموذج بوديكو؛ ويقع الجليد
قرب \ang{48} عند العتبة، مقابل \ang{50} عنده. يتجمّد نموذج خطوط العرض
أسرع من النموذج الشامل، لأنّ ضوء الشمس يكون أضعف حيث
يتكوّن الجليد.

\textbf{24.} يتجمّد الكوكب النموذجي عند انخفاض
\textbf{\qty[text-series-to-math=true]{2.9}{\percent}} من الثابت
الشمسي (\qty{2.87}{\percent} بالحساب).
\end{solution}
